%
% 4-symmetrisch.tex
%
% (c) 2025 Prof Dr Andreas Müller
%
\section{Symmetrische Polynome}
\kopfrechts{Symmetrische Polynome}%
Der Fundamentalsatz der Algebra garantiert, dass ein Polynom vom Grad
$n$ über $\mathbb{C}$ $n$ möglicherweise verschiedene Nullstellen hat.
Aus den Nullstellen lassen sich auch die Koeffizienten des Polynoms
rekonstruieren.
Dies bedeutet aber auch, dass Objekte, die als Nullstellen eines Polynoms
definiert sind, nicht eindeutig definiert sind.
Es wird daher nötig sein, zu untersuchen, wie sich daraus abgeleitete
Eigenschaften unterscheiden, wenn die Nullstellen permutiert werden.

%
% 41-elementar.tex -- Elementarsymmetrische Polynome und der Fundamentalsatz
%
% (c) 2026 Prof Dr Andreas Müller
%
\subsection{Elementarsymmetrische Polynome und der Fundamentalsatz
\label{buch:komplex:symmetrisch:subsection:fundamentalsatz}}
Der Fundamentalsatz der Algebra besagt, dass sich ein Polynom in
Linearfaktoren zerlegen lässt.
Multipliziert man die Linearfaktoren aus, entstehen die Koeffizienten
des ursprünglichen Polynoms durch algebraische Ausdrücke aus den
Nullstellen.
In diesem Abschnitt soll untersucht werden, wie diese algebraischen
Ausdrücke entstehen.

%
% Permutationen
%
\subsubsection{Permutationen}
Im Folgenden soll untersucht werden, wie sich ein Polynom von mehreren
Variablen ändert, wenn die Variablen vertauscht werden.
Wir benötigen dazu die folgenden Notationen.

\begin{definition}[Permutation]
Eine \emph{Permutation} $\sigma$ von $n$ Elementen ist eine bijektive Abbildung
\index{Permutation}%
$\pi\colon [n]\to[n]$.
Die Menge aller Permutationen von $n$ Elementen bilden die
\emph{symmetrische Gruppe} $S_n$.
\index{symmetrische Gruppe}%
\index{Gruppe!symmetrisch}%
\index{Sn@$S_n$}%
\end{definition}

Sei jetzt $f(X_1,\dots,X_n)$ eine Funktion der $n$ Variablen 
$X_1,\dots,X_n$.
Die Vertauschung der Variablen mit einer Permutation $\pi\in S_n$
ergibt eine neue Funktion $f^\pi$, die durch
\[
f^\pi(X_1,X_2,\dots,X_n)
=
f(X_{\pi(1)},X_{\pi(2)},\dots,X_{\pi(n)})
\]
definiert ist.

\begin{beispiel}
Die Funktion $f(X_1,X_2) = \cos X_1\sin X_2 - \sin X_1\cos X_2$ ist eine
Funktion von zwei Variablen.
Die Gruppe $S_2$ enthält nur zwei Elemente, das neutrale Element und
die Vertauschung $\sigma\in S_2$ der beiden Variablen.
Die Funktion
\begin{align*}
f^\sigma(X_1,X_2)
&=
f(X_{\pi(1)},X_{\pi(2)})
=
f(X_2,X_1)
\\
&=
\cos X_2\sin X_1 - \sin X_2\cos X_1
\\
&=
-\cos X_1\sin X_2 + \sin X_1\cos X_2
=
-f(X_1,X_2),
\end{align*}
es ist also $f^\sigma=-f$.
Diese Aussage könnte man auch mit dem Additionstheorem bekommen, welches
\index{Additionstheorem}%
\begin{align*}
f(X_1,X_2)&=\sin(X_1-X_2)
\\
\Rightarrow\qquad
f^\sigma(X_1,X_2)
&=
\sin(X_2-X_1)
=
-\sin(X_1-X_2)
=
-f(X_1,X_2)
\end{align*}
garantiert.
\end{beispiel}

%
% Symmetrische Polynome
%
\subsubsection{Symmetrische Polynome}
Die Reihenfolge der Linearfaktoren, in die eine Polynom zerfällt, spielt
keine Rolle.
Jede Reihenfolge führt auf das gleiche ausmultiplizierte Polynom.
Um die Entstehung der Koeffizienten zu verstehen, müssen daher nur
algebraische Ausdrücke in den Nullstellen studiert werden, die sich nicht
ändern, wenn man die Nullstellen vertauscht.
Da beim Ausmultiplizieren nur Multiplikation und Addition verwendet
werden, müssen wir also Polynome in mehreren Variablen studieren, die
sich nicht ändern, wenn man die Variablen vertauscht.

\begin{definition}[symmetrisches Polynom]
\index{symmetrisches Polynom}%
Ein Polynom $p(X_1,\dots,X_n)$ in den Variablen $X_1,\dots,X_n$
heisst \emph{symmetrisch}, wenn jede Permutation $\sigma\in S_n$
der Variablen das gleiche Polynome ergibt.
Es gilt also
\[
p(X_{\sigma(1)},\dots,X_{\sigma(n)})
=
p(X_1,\dots,X_n)
\]
für alle Permutationen $\sigma\in S_n$.
\end{definition}

\begin{beispiel}
\label{buch:komplex:fundamental:symmetrisch:bsp:n=2}
Wir bestimmen die symmetrischen Polynome für $n=2$.
Es gibt nur zwei Permutationen für $n=2$, nämlich die identische
Permutation und die Vertauschung $1\leftrightarrow 2$.

Das Polynom $p(X_1,X_2)=X_1$ ist nicht symmetrisch, da die Vertauschung
$1\leftrightarrow 2$ daraus das Polynom $X_2\ne p(X_1,X_2)$ macht.
Erst die Kombination $\sigma_1(X_1,X_2)=X_1+X_2$ ist unter der Vertauschung
invariant.

Symmetrische Polynome zweiten Grades sind zum Beispiel
$\sigma_2(X_1,X_2)=X_1X_2$, aber auch
\begin{align*}
q(X_1,X_2)
&=
X_1^2+X_2^2.
\intertext{Allerdings ist dies kein ``neues'' Polynom, denn es ist}
X_1^2+X_2^2
&=
(X_1+X_2)^2 - 2X_1X_2
\\
&=
\sigma_1(X_1,X_2)^2 - 2 \sigma_2(X_1,X_2).
\end{align*}
Das Polynom $X_1^2+X_2^2$ kann also durch $\sigma_1(X_1,X_2)$ und
$\sigma_2(X_1,X_2)$ ausgedrückt werden.

Auch für symmetrische Polynome dritten Grades können wir Beispiele
finden.
In allen Fällen lassen sich diese Polynome durch die
$\sigma_k(X_1,X_2)$, $k=1,2$, ausdrücken:
\begin{align*}
X_1^2X_2 + X_1X_2^2
&=
(X_1+X_2)X_1X_2
=
\sigma_1(X_1,X_2)\cdot\sigma_2(X_1,X_2),
\\
X_1^3 + X_2^3
&=
(X_1+X_2)^3 - 3X_1^2X_2 - 3X_1X_2^2
\\
&=
(X_1+X_2)^3 - 3(X_1^2X_2 + X_1X_2^2)
\\
&=
\sigma_1(X_1,X_2)^3 - 3\sigma_1(X_1,X_2)\cdot \sigma_2(X_1,X_2).
\intertext{Dasselbe gilt für symmetrische Polynome vom Grad 4:}
X_1^2 X_2^2
&=
\sigma_2(X_1,X_2)^2
\\
X_1^3X_2+X_1X_2^3
&=
X_1X_2(X_1^2+X_2^2)
\\
&=
\sigma_2(X_1,X_2)\bigl(\sigma_1(X_1,X_2)^2-2\sigma_2(X_1,X_2)\bigr)
\\
&=
\sigma_1(X_1,X_2)^2 \sigma_2(X_1,X_2)
-2
\sigma_2(X_1,X_2)^2
\\
X_1^4+X_2^4
&=
\sigma_1(X_1,X_2)^4
-
4\sigma_1(X_1,X_2)^2 \sigma_2(X_1,X_2)
+8
\sigma_2(X_1,X_2)^2
-
6\sigma_2(X_1,X_2)^2
\\
&=
\sigma_1(X_1,X_2)^4
-
4\sigma_1(X_1,X_2)^2 \sigma_2(X_1,X_2)
+2
\sigma_2(X_1,X_2)^2.
\end{align*}
Keines der symmetrischen Polynome dritten oder vierten Grades enthält
Information, die sich nicht auch schon aus den beiden speziellen
Polynomen $\sigma_1(X_1,X_2)$ und $\sigma_2(X_1,X_2)$ ablesen lässt.
\end{beispiel}

Im Beispiel wurden zwei Polynome gefunden, mit denen sich alle 
anderen symmetrischen Polynome ausdrücken liessen.
Dies sind die sogenannten elemetarsymmetrischen Polynome,
die im nächsten Abschnitt eingeführt werden.

%
% Elementarsymmetrische Polynome
%
\subsubsection{Elementarsymmetrische Polynome}
Wir betrachten das Polynom
\[
p(T,X_1,\dots,X_n)
=
(T+X_1)\cdot\ldots\cdot (T+X_n)
\]
in den Variablen $T,X_1,\dots,X_n$.
Da das Produkt kommutativ ist, ist dieses Polynom symmetrisch in den
Variablen $X_1,\dots,X_n$.
Multipliziert man das Polynom aus und fasst die Terme mit gleicher
Potenz von $T$ zusammen erhält man
\begin{equation}
\renewcommand{\arraycolsep}{2pt}
\begin{array}{rclclclcl}
p(T,X_1,\dots,X_n)
&=&
T^n
&+&
\sigma_1(X_1,\dots,X_n)
T^{n-1}
&+&
\sigma_2(X_1,\dots,X_n)
T^{n-2}
&+&
\ldots
\\
& &
\ldots
&+&
\sigma_{n-1}(X_1,\dots,X_n) T
&+&
\sigma_n(X_1,\dots,X_n)
\\
&=&
\rlap{$\displaystyle \sum_{k=0}^n \sigma_k(X_1,\dots,X_n) T^{n-k} $}
\end{array}
\label{buch:komplex:symmetrisch:eqn:definition}
\end{equation}
mit Polynomen $\sigma_k(X_1,\dots,X_n)$, $k=0,\dots,n$.
Da $p(T,X_1,\dots,X_n)$ symmetrisch ist, gilt
\[
p(T,X_{\sigma(1)},\dots,X_{\sigma(n)})
=
p(T,X_1,\dots,X_n)
\] 
für jede Permutation $\sigma\in S_n$.
Durch Koeffizientenvergleich folgt, dass
\[
\sigma_k(X_{\sigma(1)},\dots,X_{\sigma(n)})
=
\sigma_k(X_1,\dots,X_n)
\]
für jeden Koeffizienten $\sigma_k$, $k=1,\dots,n$, und für jede
Permutation $\sigma\in S_n$.
Insbesondere sind die $\sigma_k(X_1,\dots,X_n)$ alle symmetrische
Polynome.

\begin{definition}[Elementarsymmetrische Polynome]
\label{buch:komplex:fundamental:def:elementarsymmetrisch}
Die durch \eqref{buch:komplex:symmetrisch:eqn:definition} definierten
Polynome $\sigma_k(X_1,\dots,X_n)\in \mathbb{Z}[X_1,\dots,X_n]$ 
heissen die \emph{elementarsymmetrischen Polynome}.
\index{elementarsymmetrische Polynome}%
\end{definition}

\begin{beispiel}
Für $n=2$ entstehen die elementarsymmetrischen Polynome aus dem
Produkt
\[
(T+X_1)(T+X_2)
=
T^2 +(X_1+X_2)T+X_1X_2
\quad\Rightarrow\quad
\left\{
\begin{aligned}
\sigma_1(X_1,X_2) &= X_1+X_2 \\
\sigma_2(X_1,X_2) &= X_1X_2.
\end{aligned}
\right.
\]
Dies sind genau die Polynome, die im
Beispiel~\ref{buch:komplex:fundamental:symmetrisch:bsp:n=2}
gefunden wurden.
\end{beispiel}

\begin{beispiel}
Für $n=3$ 
kann man aus dem Produkt
\[
(T+X_1)
(T+X_2)
(T+X_3)
=
T^3
+
(X_1+X_2+X_3)T^2
+(X_1X_2+X_1X_3+X_2X_3)T
+
X_1X_2X_3
\]
die elementarsymmetrischen Polynome
\begin{equation}
\begin{aligned}
\sigma_1(X_1,X_2,X_3) &= X_1+X_2+X_3\\
\sigma_2(X_1,X_2,X_3) &= X_1X_2 + X_1X_3 + X_2X_3\\
\sigma_3(X_1,X_2,X_3) &= X_1X_2X_3
\end{aligned}
\label{buch:komplex:fundamental:bsp:elementar3}
\end{equation}
ablesen.
\end{beispiel}

Die Beispiele zeigen auch, dass für beliebiges $n$ die
elementarsymmetrischen Polynome für $k=0$, $k=1$ und 
$k=n$
die besonders einfache Form
\begin{align*}
\sigma_0(X_1,\dots,X_n) &= 1 \\
\sigma_1(X_1,\dots,X_n) &= X_1+\ldots + X_n \\
\sigma_n(X_1,\dots,X_n) &= X_1 \cdots X_n
\end{align*}
haben.
Die anderen elementarsymmetrischen Polynome können als
\[
\sigma_k(X_1,\dots,X_n)
=
\sum_{1\le i_1<i_2<\dots<i_k\le n}
X_{i_1}
X_{i_2}
\cdots
X_{i_k}
\]
geschrieben werden.

Die elementarsymmetrischen Polynome können dazu verwendet werden, die
Koeffizienten eines Polynoms aus den Nullstellen zu berechnen.
Die Nullstellen eines solchen Polynoms sind algebraische Zahlen über
dem Koeffizientenkörper.
Oft sind alle Nullstellen ausserhalb des Koeffzientenkörpers.
Trotzdem sind die Werte der elementarsymmetrischen Funktionen, ausgewertet
auf den Nullstellen des Polynoms, wieder im Koeffizientenkörper.

\begin{beispiel}
Das Polynom $p(X) = X^2+pX+q$ hat die Nullstellen
\[
\alpha_{\pm}
=
-\frac{p}{2}
\pm
\!
\sqrt{ \frac{p^2}{4}-q }.
\]
Setzt man diese Nullstellen in die elementarsymmetrischen Funktionen
für $n=2$ ein, erhält man
\[
\renewcommand{\arraycolsep}{1.5pt}
\begin{array}{rccclcl}
\sigma_1(\alpha_+,\alpha_-)
&=&
\alpha_++\alpha_-
&=&
\displaystyle
-\frac{p}{2}
+
\!\sqrt{ \frac{p^2}{4}-q }
+
-\frac{p}{2}
-
\!\sqrt{ \frac{p^2}{4}-q }
&=&
-p,
\\[7pt]
\sigma_2(\alpha_+,\alpha_-)
&=&
\alpha_+\cdot \alpha_-
&=&
\displaystyle
\biggl(
-\frac{p}{2}
+
\!\sqrt{ \frac{p^2}{4}-q }
\biggr)
\biggl(
-\frac{p}{2}
-
\!\sqrt{ \frac{p^2}{4}-q }
\biggr)
\\[7pt]
& &
&=&
\displaystyle
\biggl(-\frac{p}{2}\biggr)^2
-
\biggl(
\!\sqrt{ \frac{p^2}{4}-q }
\biggr)^2
\\[7pt]
& &
&=&
\displaystyle
\frac{p^2}{4}
-
\frac{p^2}{4}+q
&=&
q.
\end{array}
\]
Die Werte der elementarsymmetrischen Funktionen von $\alpha_{\pm}$ sind
also wieder im Koeffizientenkörper.
\end{beispiel}

\begin{satz}
Sei $p(X) = p_0+p_1X+\dots +p_{n-1}X^{n-1} + X^n\in K[X]$ ein normiertes
mit Koeffizienten in einem Körper $K$ und seien $\alpha_1,\dots,\alpha_n$
alle Nullstellen, so dass 
\[
p(X)
=
(X-\alpha_1)(X-\alpha_2)\cdots(X-\alpha_n) =
\prod_{k=1}^n (X-\alpha_k).
\]
Dann ist 
\[
(-1)^k
\sigma_k(\alpha_1,\dots,\alpha_n) 
=
p_k
\]
für $k=1,\dots,n$ oder
\[
p(X)
=
\sum_{k=0}^n
(-1)^k
\sigma_k(\alpha_1,\dots,\alpha_n) X^{n-k}.
\]
Ausserdem ist
\[
\sigma_k(\alpha_1,\dots,\alpha_n) \in K
\]
für $k=1,\dots,n$.
\end{satz}

\begin{beispiel}
Für das Polynom $p(X) = X^5 + 2X^4+3X^3+4X^2+5X+6\in\mathbb{Z}[X]$
kann man numerisch (zum Beispiel mit Maxima) die Nullstellen
\[
\renewcommand{\arraycolsep}{1.5pt}
\begin{array}{rcrcl}
\alpha_1 &=&  0.5516854634589816 &+& 1.2533488602772058 i\\
\alpha_2 &=&  0.5516854634589816 &-& 1.2533488602772058 i\\
\alpha_3 &=& -0.8057864693890309 &+& 1.2229047133744104 i\\
\alpha_4 &=& -0.8057864693890309 &-& 1.2229047133744104 i\\
\alpha_5 &=& -1.4917979881399015 & &
\end{array}
\]
finden.
Die Werte der elementarsymmetrischen Polynome müssen wieder
ganze Zahlen ergeben. 
Das Polynom $\sigma_1(X_1,\dots,X_5)$ ist die Summe
und $\sigma_5(X_1,\dots,X_5)$ ist das Produkt der Argumente.
Tatsächlich sind Summe und Produkt
\[
\renewcommand{\arraycolsep}{1.5pt}
\begin{array}{rcccl}
\sigma_1(\alpha_1,\dots,\alpha_5)
&=&
\displaystyle
\sum_{k=1}^5 \alpha_k
&=&
-2
\\
\sigma_5(\alpha_1,\dots,\alpha_5)
&=&
\displaystyle
\prod_{k=1}^5 \alpha_k
&=&
-6.000000000000004 + 3.525489416408166\cdot 10^{-16}i
\end{array}
\]
mit Maschinengenauigkeit bis auf das Vorzeichen die Koeffizienten
des Polynoms $p(X)$.
\end{beispiel}

%
% Der Fundamentalsatz für symmetrische Polynome
%
\subsubsection{Der Fundamentalsatz für symmetrische Polynome}
Die elementarsymmetrischen Polynome sind in einem gewissen Sinne
die einzigen symmetrischen Polynome, die man braucht.
Dies ist der Inhalt des folgenden Fundamentalsatzes für
symmetrische Polynome: jedes symmetrische Polynom lässt sich
als Polynom in den elementarsymmetrischen Polynomen ausdrücken.

\begin{satz}[Fundamentalsatz für symmetrische Polynome]
\label{buch:komplex:fundamental:satz:fundamentalpolynom}
Ist $p(X_1,\dots,X_n)\in K[X_1,\dots,X_n]$ ein symmetrisches
Polynom von $n$ Variablen, dann gibt es genau ein Polynom
$q(S_1,\dots,S_n)\in K[S_1,\dots,S_n]$ derart, dass
\index{Fundamentalsatz für symmetrische Polynome}%
\[
p(X_1,\dots,X_n)
=
q\bigl(\sigma_1(X_1,\dots,X_n),\dots,\sigma_n(X_1,\dots,X_n)\bigr).
\qedhere
\]
\end{satz}

Nicht nur gibt es das Polynom $q$, wir werden einen Beweis
geben, der das Polynom sogar rekursiv konstruiert.
Wir illustrieren die Idee am Beispiel eines symmetrischen
Polynoms $q(X_1,X_2)$ dritten Grades in $X_1$ und $X_2$.
Es muss mindestens einen der Terme
\[
X_1^3,\;
X_1^2X_2,\;
X_1X_2^2,\;
X_2^3
\]
dritten Grades enthalten.
Falls der Term $cX_1^3$ in $q$ vorkommt, dann muss auch der Term $cX_2^3$
vorkommen.
Mit dem elementarsymmetrischen Polynom $\sigma_1(X_1,X_2)$ können
wir ebenfalls ein Polynom dritten Grades konstruieren, welches
$cX_1^3$ und $cX_2^3$ enthält, nämlich $c\sigma_1(X_1,X_2)^3$.
Die Differenz $q_1(X_1,X_2)=q(X_1,X_2)-c\sigma_1(X_1,X_2)^3$ enthält die
Terme $X_1^3$ und $X_2^3$ nicht mehr.

$q_1(X_1,X_2)$ kann aber immer noch $c_1X_1^2X_2$ enthalten.
Symmetrie verlangt dann, dass es auch $c_1X_1X_2^2$ enthält.
Mit den elementarsymmetrischen Polynomen können wir 
\[
c_1\sigma_1(X_1,X_2)\sigma_2(X_1,X_2)
=
c_1(X_1+X_2)(X_1X_2)
=
c_1(X_1^2X_2+X_1X_2^2)
\]
konstruieren, welches die gleichen Terme enthält.
Das neue Polynom
$q_2(X_1,X_2) = q_1(X_1,X_2) - c_1\sigma_1(X_1,X_2)\sigma_2(X_1,X_2)$
enthält dann auch diese Terme nicht mehr.

Da es keine weiteren Terme dritten Grades gibt, ist $q_2$ jetzt nur noch
ein Polynom zweiten Grades.
Falls es den Term $c_2X_1^2$ enthält, können wir diesen durch
$q_3(X_1,X_2)=q_2(X_1,X_2)-c_2\sigma_1(X_1,X_2)^2$ eliminieren.
Falls dann noch ein Term $c_3X_1X_2$ bleibt, wird dieser durch
$q_4(X_1,X_2)=q_3(X_1,X_2)-c_3\sigma_2(X_1,X_2)$ eliminiert.

$q_4(X_1,X_2)$ ist ein Polynom ersten Grades und kann daher sofort
als Linearkombination von $\sigma_1(X_1,X_2)=X_1+X_2$ und
$\sigma_0(X_1,X_2)=1$ dargestellt werden.

Der Algorithmus eliminiert also die verschiedenen Arten von Monomen, die
in $q(X_1,X_2)$ vorkommen können dadurch, dass ein geeignetes Produkt von 
elementarsymmetrischen Polynomen subtrahiert wird.
Dabei ist in einer Reihenfolge vorzugehen, die verhindert, dass bereits
eliminierte Terme wieder hinzugefügt werden.
Es muss also als vor allem eine geeignete Anordnung der Monome definiert
werden.

Um die Ordnung zweier Monome gleichen Grades zu definieren, kürzen wir 
das Monom $X_1^{a_1}\dots X_n^{a_n}$ durch das $n$-Tupel
\[
e(X_1^{a_1}\cdots X_n^{a_n})
=
(a_n,a_{n-1},\dots,a_2,a_1)
\]
ab.
Die Funktion $e$ extrahiert also die Exponenten der einzelnen 
Variablen und ordnet sie nach absteigendem Index der Variablen an.
Zum $n$-Tupel $(a_n,\dots,a_1)$ ist $m=X_1^{a_1}\cdots X_n^{a_n}$
das Monom mit $e(m)=(a_n,\dots,a_1)$.

\begin{definition}[umgekehrt lexikographische Ordnung]
\label{buch:komplex:symmetrisch:def:ulexiko}
\index{umgekehrt lexikographische Ordnung}%
Die umgekehrt lexikographische Ordnung von $n$-Tupeln ist definiert durch
\[
(a_n,\dots,a_1) < (b_n,\dots,b_1)
\quad
:\Leftrightarrow
\quad
\renewcommand{\arraycolsep}{2pt}
\left\{\begin{array}{rclcl}
a_n &>& b_n	&\quad&\text{falls $a_n\ne b_n$}\\
a_{n-1} &>& b_{n-1}&&\text{falls $a_{n-1}\ne b_{n-1}$ und $a_n=b_n$}\\[-2pt]
&\vdots&        &&\vdots\\
a_2 &>& b_2	&&\text{falls $a_2\ne b_2$ und $a_k=b_k$ für $k\ge 3$}\\
a_1 &>& b_1	&&\text{falls $a_1\ne b_1$ und $a_k=b_k$ für $k\ge 2$.}
\end{array}
\right.
\]
Zwei Monome $m_1$ und $m_2$ gleichen Grades erfüllen $m_1 < m_2$, wenn
$e(m_1) < e(m_2)$.
\end{definition}

Die umgekehrt lexikographische Ordnung der Tripel vom Grad 3 ist
\index{lexikographisch}%
\begin{equation}
\begin{aligned}
(3,0,0)
&<
(2,1,0)
<
(2,0,1)
<
(1,2,0)
<
(1,1,1)
\\
&<
(1,0,2)
<
(0,3,0)
<
(0,2,1)
<
(0,1,2)
<
(0,0,3).
\end{aligned}
\label{buch:komplex:symmetrisch:eqn:3tupelordnung}
\end{equation}
``Grosse'' Elemente im Sinn der umgekehrt lexikographischen Ordnung
haben also möglichst grosse Komponenten $a_i$ mit kleinem Index $i$.

\begin{lemma}
\label{buch:komplex:symmetrisch:lemma:maximal}
Ein $n$-Tupel $(a_n,\dots,a_1)$ ist genau dann das grösste im Sinne
der umgekehrt lexikographischen Ordnung, wenn
$a_n\le a_{n-1}\le \ldots \le a_1$
ist.
\end{lemma}

\begin{proof}
Wäre $a_{i+1} > a_i$, dann liesse sich ein grösseres $n$-Tupel durch
Vertauschen von $a_i$ und $a_{i+1}$ finden:
\[
(a_n,\dots,a_{i+1},a_i,\dots,a_1)
<
(a_n,\dots,a_i,a_{i+1},\dots,a_1).
\]
Dies zeigt, dass das $n$-Tupel $(a_n,\dots,a_1)$ nicht das grösste sein 
kann, wenn die $a_k$ nicht monoton sind.
\end{proof}

Ein $n$-Tupel wie in Lemma~\ref{buch:komplex:symmetrisch:lemma:maximal}
wird auch ein \emph{maximales $n$-Tupel} genannt.
\index{maximales n-Tupel@maximales $n$-Tupel}%

\begin{definition}[Monomordnung]
\index{Monomordnung}%
Für zwei beliebige Monome $m_1$ und $m_2$ von möglicherweise verschiedenem
Grad wird die Ordnung durch
\[
m_1 < m_2
\quad
:\Leftrightarrow
\quad
\begin{cases}
\deg m_1 < \deg m_2&\qquad\text{falls $\deg m_1\ne \deg m_2$}\\
m_1 < m_2&\qquad\text{falls $\deg m_1 = \deg m_2$}
\end{cases}
\]
definiert, wobei die Ordnung gleichgradiger Monome im zweiten Fall
die ungekehrt lexikographische Ordnung von
Definition~\ref{buch:komplex:symmetrisch:def:ulexiko} ist.
\end{definition}

Für zwei Variablen $X_1$ und $X_2$ ergeben sich bis zum Grad 3 die folgenden 
Reihenfolgen:
\begin{align}
&\text{Grad 1:} & &X_1 < X_2
\notag
\\
&\text{Grad 2:} & &X_1^2 < X_1X_2 < X_2^2
\notag
\\
&\text{Grad 3:} & &
X_3^3
<
X_3^2 X_2
<
X_3^2 X_1
<
X_3 X_2^2
<
X_3 X_2 X_1
<
X_3 X_1^2
<
X_2^3
<
X_2^2
X_1
<
X_2
X_1^2
<
X_1^3
\label{buch:komplex:symmetrisch:eqn:3monomordung}
\end{align}
wobei sich
\eqref{buch:komplex:symmetrisch:eqn:3monomordung}
aus
\eqref{buch:komplex:symmetrisch:eqn:3tupelordnung}
ergibt.

\begin{definition}[führendes Monom, führender Koeffizient]
Das \emph{führende Monom} oder \emph{leading monimial}
$\operatorname{LM}(p)$ eines Polynoms $p$ ist
\index{fuhrendes Monom@führendes Monom}%
\index{leading monomial}%
\index{LM}%
\index{Monom!fuhrend@Monom!führend}%
das grösste Monom, welches in $p$ vorkommt.
Der \emph{führende Koeffizient} oder \emph{leading coefficient}
$\operatorname{LC}(p)$ von $p$ 
ist der Koeffizient des führenden Monoms in $p$.
\index{fuhrendes Koeffizient@führendes Koeffizient}%
\index{leading coefficient}%
\index{LC}%
\index{Koeffizient!fuhrend@Koeffizient!führend}%
\end{definition}

\begin{lemma}
Für $k=1,\dots,n$ gilt
\[
\operatorname{LM}(\sigma_k(X_1,\dots,X_n))
=
X_1\cdots X_k
\quad\text{mit}\quad
e(\operatorname{LM}(\sigma_k(X_1,\dots,X_n)))
=
(\underbrace{0,\dots,0}_{n-k},\underbrace{1,\dots,1}_{k}).
\]
\end{lemma}

\begin{proof}
Die Monome eines elementarsymmetrischen Polynoms enthalten keine
höheren als erste Potenzen jeder Variable.
In der definierten Ordnung hat das führend Monom daher nur die $k$
``grössten'' Variablen.
\end{proof}

\begin{lemma}
Ist $(a_n,\dots,a_1)$ ein maximales $n$-Tupel, dann ist
\[
m
=
\operatorname{LM}(
\sigma_1^{a_1-a_2}
\sigma_2^{a_2-a_3}
\cdots
\sigma_{n-1}^{a_{n-1}-a_n}
\sigma_{n}^{a_n}
)
\]
ein Monom mit $e(m)=(a_n,\dots,a_1)$.
\end{lemma}

\begin{proof}
Statt mit den Monomen zu rechnen, können wir mit den $n$-Tupeln rechnen:
\begin{align*}
\left.
\renewcommand{\arraycolsep}{2pt}
\begin{array}{crcl}
 & (0,\dots,0,0,1) &\cdot & (a_1-a_2)     \\
+& (0,\dots,0,1,1) &\cdot & (a_2-a_3)     \\
+& (0,\dots,1,1,1) &\cdot & (a_3-a_4)     \\[-2pt]
 &		   &\vdots&               \\
+& (0,\dots,1,1,1) &\cdot & (a_{n-1}-a_n) \\
+& (1,\dots,1,1,1) &\cdot & \phantom{(}a_n\\
\end{array}
\right\}
&=
\left\{
\renewcommand{\arraycolsep}{1pt}
%\begin{array}{clccclclcll}
\begin{array}{ccccccccccl}
 &(0       &,& \dots &,& 0           &,&       0     &,& a_1-a_2     &)\\
 &(0       &,& \dots &,& 0           &,& a_2-a_3     &,& a_2-a_3     &)\\
 &(0       &,& \dots &,& a_3-a_2     &,& a_3-a_4     &,& a_3-a_4     &)\\[-2pt]
 &\;\vdots & & \vdots& & \vdots      & & \vdots      & & \vdots      & \\
 &(0       &,& \dots &,& a_{n-1}-a_n &,& a_{n-1}-a_n &,& a_{n-1}-a_n &)\\
 &(a_n     &,& \dots &,& a_n         &,&  a_n        &,& a_n         &)
\end{array}
\right.
\\
&=
(a_n,\dots,a_3,a_2,a_1),
\end{align*}
wie verlangt.
\end{proof}

Nach diesen Vorbereitungen sind wir jetzt in der Lage, den
Satz~\ref{buch:komplex:fundamental:satz:fundamentalpolynom}
dadurch zu beweisen, dass wir den in den Vorbemerkungen zum Beweis
skizzierten Algorithmus ausformulieren.

\begin{proof}[Beweis von Satz~\ref{buch:komplex:fundamental:satz:fundamentalpolynom}]
Das Polynom $q$ wird iterativ durch den folgenden Algorithmus bestimmt:
\begin{enumerate}
\item Setze $q=0$.
\item Setze $a=e(\operatorname{LM}(p))$.
\item Ersetze $p$ und $q$ durch
\begin{align*}
p&\leftarrow
p - \operatorname{LC}(p) 
\sigma_1^{a_2-a_1}\sigma_2^{a_3-a_1}\cdots\sigma_{n-1}^{a_n-a_{n-1}},
\\
q&\leftarrow
q+ \operatorname{LC}(p) 
\sigma_1^{a_2-a_1}\sigma_2^{a_3-a_1}\cdots\sigma_{n-1}^{a_n-a_{n-1}}.
\end{align*}
\item Falls $p\ne 0$ weiter bei 2.
\end{enumerate}
Im Schritt~3 wird $p$ durch ein Polynom mit kleinerem führendem 
Monom ersetzt.
Nach endlich vielen Schritten terminiert der Algorithmus daher und
$q$ ist das gesuchte Polynom in den elementarsymmetrischen Funktionen.
\end{proof}

\begin{beispiel}
Wir führen den Algorithmus im Fall $n=3$ für das Polynom
\[
p(X_1,X_2,X_3) = {\color{darkred}X_1^3} + X_2^3 + X_3^3
\]
durch.
\smallskip

\noindent
Das führende Monom von $p$ ist 
\(
\operatorname{LM}(p) = {\color{darkred}X_1^3}
\)
mit führendem Koeffizienten $\operatorname{LC}(p)=1$,
somit sind die neuen Polynome $p$ und $q$
\begin{align*}
p
&\leftarrow
p-\sigma_1^3
= 
- 6 X_3 X_2 X_1
- 3(
  X_3^2 X_2
+ X_3^2 X_1
+ X_3 X_2^2
+ X_3 X_1^2
+ X_2^2 X_1
+ {\color{darkred}X_2 X_1^2}
)
\\
q
&\leftarrow
\sigma_1^3.
\intertext{Für das neue Polynom ist das führende Monom das rot
hervorgehobene $\operatorname{LM}(p)=X_2X_1^2$, mit
führendem Koeffizienten $\operatorname{LC}(p)=-3$, es muss als im nächsten
Schritt $3\sigma_1\sigma_2$ addiert werden:}
p
&\leftarrow
p+3\sigma_1\sigma_2
=
3 {\color{darkred}X_3X_2X_1}
\\
q
&\leftarrow
\sigma_1^3-3\sigma_1\sigma_2.
\intertext{Das führende Monom von $p$ ist jetzt
$\operatorname{LM}(p)={\color{darkred}X_3X_2X_1}=\sigma_3$ mit
führendem Koeffizienten $\operatorname{LC}(p)=3$,
sodass im nächsten Schritt $3\sigma_3$
subtrahiert werden muss:}
p
&\leftarrow
p - 3\sigma_3 = 0
\\
q
&\leftarrow
\sigma_1^3-3\sigma_1\sigma_2+3\sigma_3.
\end{align*}
Damit haben wir
\begin{equation}
X_1^3+X_2^3+X_3^3
=
\sigma_1(X_1,X_2,X_3)^3
-
3\sigma_1(X_1,X_2,X_3)\sigma_2(X_1,X_2,X_3)
+
3\sigma_3(X_1,X_2,X_3)
\label{buch:komplex:symmetrisch:eqn:X33}
\end{equation}
gefunden.
\end{beispiel}

%
% 42-matrix.tex -- Matrixdarstellung algebraischer Elemente
%
% (c) 2026 Prof Dr Andreas Müller
%
\subsection{Matrixdarstellung algebraischer Elemente}
Für die Rechnung in einem Erweiterungskörper $\mathbb{Q}(\alpha)$
mit dem Minimalpolynom $m(X)\in \mathbb{Q}[X]$ ist bis jetzt kein
einfacher Algorithmus gegeben worden.
Die Matrixdarstellung eines algebraischen Elementes, die in diesem
Abschnitt studiert werden soll, liefert so einen Algorithmus.
Gleichzeitig wird damit aber auch ein Werkzeug bereitgestellt,
welches beim Beweis des Liouville-Prinzips für algebraische Erweiterungen
\index{Liouville-Prinzip}%
im Abschnitt~\ref{buch:intalg:liouville:subsection:allgemein}
als nützlich erweisen wird.

%
% Die Matrix eines algebraischen Elementes
%
\subsubsection{Die Matrix eines algebraischen Elements}
In diesem Abschnitt betrachten wir die Multiplikationsoperation
mit $\alpha$ im Erweiterungskörper $\mathbb{Q}(\alpha)$ für ein
algebraisches Element mit dem Minimalpolynom
\[
m(X)
=
m_0+m_1X+\dots + m_{n-1}X^{n-1}+X^n
\in
\mathbb{Q}[X].
\]
Der Erweiterungskörper ist die Menge
\[
\mathbb{Q}(\alpha)
=
\{
a_0+a_1\alpha+\ldots+a_{n-1}\alpha^{n-1}
\mid
a_k\in\mathbb{Q}\;
\forall k
\}.
\]
Er ist ein $n$-dimensionaler Vektorraum über $\mathbb{Q}$ mit der
Basis
\[
\mathscr{B}
=
\{
1,\alpha,\alpha^2,\ldots,\alpha^{n-1}
\}.
\]
Die Multiplikation mit $\alpha$ ist wegen
\[
\alpha (t u + s v)
=
t
\alpha
u
+
s
\alpha
v
,\qquad t,s\in\mathbb{Q},\;u,v\in\mathbb{Q}(\alpha)
\]
eine lineare Abbildung, man kann sie also auch in der Basis
$\mathscr{B}$ durch eine $n\times n$-Matrix $A$ beschreiben.
Aus der Abbildung der Basiselemente in $\mathscr{B}$ kann
man die Matrix
\begin{equation}
\left.
\begin{aligned}
1           &\mapsto \alpha   \\
\alpha      &\mapsto \alpha^2 \\
\alpha^2    &\mapsto \alpha^3 \\
\alpha^3    &\mapsto \alpha^4 \\[-3pt]
            &\phantom{n}\vdots\\
\alpha^{n-1}&\mapsto -a_0-a_1\alpha-\dots-a_{n-1}\alpha^{n-1}
\end{aligned}
\right\}
\quad\Rightarrow\quad
\bm{\alpha}
=
\left(
\raisebox{-1.39cm}{\begin{tikzpicture}[>=latex,thick]
\def\punkt#1#2{ ({(#1)*0.75},{-(#2)*0.5}) }
\clip \punkt{-0.5}{-0.5} rectangle \punkt{5.5}{5.5};
%\draw[color=darkred] \punkt{-0.5}{-0.5} rectangle \punkt{5.5}{5.5};
\def\element#1#2#3{
\node at \punkt{#1}{#2} {$#3\mathstrut$};
}
\def\horizontaldots#1#2{
\fill \punkt{#1-0.333}{#2} circle[radius=0.02];
\fill \punkt{#1}{#2} circle[radius=0.02];
\fill \punkt{#1+0.333}{#2} circle[radius=0.02];
}
\def\verticaldots#1#2{
\fill \punkt{#1}{#2-0.333} circle[radius=0.02];
\fill \punkt{#1}{#2} circle[radius=0.02];
\fill \punkt{#1}{#2+0.333} circle[radius=0.02];
}
\def\diagonaldots#1#2{
\fill \punkt{#1-0.333}{#2-0.333} circle[radius=0.02];
\fill \punkt{#1}{#2} circle[radius=0.02];
\fill \punkt{#1+0.333}{#2+0.333} circle[radius=0.02];
}
\draw[color=darkred!10,line width=11pt] \punkt{-1}{-1} -- \punkt{6}{6};
\draw[color=blue!10,line width=11pt] \punkt{-1}{0} -- \punkt{5}{6};
\element{0}{0}{0} \element{1}{0}{0} \element{2}{0}{0}
\horizontaldots{3}{0}
\element{4}{0}{0} \element{5}{0}{-m_0}
%
\element{0}{1}{1} \element{1}{1}{0} \element{2}{1}{0}
\horizontaldots{3}{1}
\element{4}{1}{0} \element{5}{1}{-m_1}
%
\element{0}{2}{0} \element{1}{2}{1} \element{2}{2}{0}
\horizontaldots{3}{2}
\element{4}{2}{0} \element{5}{2}{-m_2}
%
\element{0}{3}{0} \element{1}{3}{0} \element{2}{3}{1}
\diagonaldots{3}{3}
\element{5}{3}{-m_3}
%
\verticaldots{0}{4}
\verticaldots{1}{4}
\verticaldots{2}{4}
\diagonaldots{3}{4}
\diagonaldots{4}{4}
\verticaldots{5}{4}
%
\element{0}{5}{0} \element{1}{5}{0} \element{2}{5}{0}
\horizontaldots{3}{5}
\element{4}{5}{1} \element{5}{5}{-m_{n-1}}
\end{tikzpicture}}
\right)
\label{buch:komplex:fundamental:eqn:alphaMatrix}
\end{equation}
ablesen.

%
% Multiplikative Inverse in Q(\alpha)
%
\subsubsection{Multiplikative Inverse in $\mathbb{Q}(\alpha)$}
Ein Element $x$ von $\mathbb{Q}(\alpha)$ schreiben wir als Spaltenvektor mit
den Einträgen $x_0,x_1,\dots,x_{n-1}$.
Die Multiplikation mit $\alpha$ wird durch die Matrix $\bm{\alpha}$ von
\eqref{buch:komplex:fundamental:eqn:alphaMatrix} beschrieben.
Die Multiplikation des Elements $a\in \mathbb{Q}(\alpha)$ mit $x$ wird
daher durch die Matrixmultiplikation mit der Matrix
\[
A
=
a_0I+a_1\bm{\alpha} + a_2\bm{\alpha}^2 + \ldots + a_{n-1}\bm{\alpha}^{n-1}
\]
wiedergegeben.
Um die Inverse $x=a^{-1}$ zu finden, muss man also das lineare
Gleichungssystem
\index{lineares Gleichungssystem}%
$ Ax=e_1 $ lösen, wobei $e_1$ der erste Standardbasisvektor ist.
\index{Standardbasisvektor}%
Man kann also das multiplikative inverse Element von $a$ als
\index{multiplikative Inverse}%
\index{Inverse, multiplikativ}%
\[
x = A^{-1}e_1
\]
finden.

%
% Das charakteristische Polynom der Matrix A
%
\subsubsection{Das charakteristische Polynom der Matrix $\bm{\alpha}$}
Das charakteristische Polynom von $\bm{\alpha}$ kann am einfachsten
durch Entwicklung nach der letzten Spalte bestimmt werden.
\index{Entwicklung}%
Die Minordeterminante zum Element $-m_k$ mit $k<n-1$ ist
\index{Minordeterminante}%
\[
\det \bm{\alpha}_{k+1,n}
=
\left|
\;
\raisebox{-2.2cm}{\begin{tikzpicture}[>=latex,thick]
\def\punkt#1#2{ ({(#1)*0.75},{-(#2)*0.5}) }
\clip \punkt{-0.5}{-0.5} rectangle \punkt{8.5}{8.5};
\begin{scope}
	\clip \punkt{-0.5}{-0.5} rectangle \punkt{4.5}{4.5};
	\draw[color=darkred!10,line width=12pt] \punkt{-1}{-1} -- \punkt{4.5}{4.5};
	\draw[color=blue!10,line width=12pt] \punkt{-1}{0} -- \punkt{4.5}{5.5};
\end{scope}
\begin{scope}
	\clip \punkt{9.5}{9.5} rectangle \punkt{4.5}{4.5};
	\draw[color=blue!10,line width=12pt] \punkt{9}{9} -- \punkt{4.5}{4.5};
	\draw[color=darkred!10,line width=12pt] \punkt{9}{8} -- \punkt{4.5}{3.5};
\end{scope}
\draw[line width=0.4pt] \punkt{-0.5}{4.5} -- \punkt{8.5}{4.5};
\draw[line width=0.4pt] \punkt{4.5}{-0.5} -- \punkt{4.5}{8.5};

%\draw[color=darkred] \punkt{-0.5}{-0.5} rectangle \punkt{8.5}{8.5};
\def\element#1#2#3{
\node at \punkt{#1}{#2} {$#3\mathstrut$};
}
\def\horizontaldots#1#2{
\fill \punkt{#1-0.333}{#2} circle[radius=0.02];
\fill \punkt{#1}{#2} circle[radius=0.02];
\fill \punkt{#1+0.333}{#2} circle[radius=0.02];
}
\def\verticaldots#1#2{
\fill \punkt{#1}{#2-0.333} circle[radius=0.02];
\fill \punkt{#1}{#2} circle[radius=0.02];
\fill \punkt{#1}{#2+0.333} circle[radius=0.02];
}
\def\diagonaldots#1#2{
\fill \punkt{#1-0.333}{#2-0.333} circle[radius=0.02];
\fill \punkt{#1}{#2} circle[radius=0.02];
\fill \punkt{#1+0.333}{#2+0.333} circle[radius=0.02];
}
\element{0}{0}{-\lambda} \element{1}{0}{0}        \element{2}{0}{0}
\horizontaldots{3}{0}
\element{4}{0}{0}        \element{5}{0}{0}        \element{6}{0}{0}
\horizontaldots{7}{0}
\element{8}{0}{0}
%
\element{0}{1}{1}        \element{1}{1}{-\lambda} \element{2}{1}{0}
\horizontaldots{3}{1}
\element{4}{1}{0}        \element{5}{1}{0}        \element{6}{1}{0}
\horizontaldots{7}{1}
\element{8}{1}{0}
%
\element{0}{2}{0}        \element{1}{2}{1}        \element{2}{2}{-\lambda}
\horizontaldots{3}{2}
\element{4}{2}{0}        \element{5}{2}{0}        \element{6}{2}{0}
\horizontaldots{7}{2}
\element{8}{2}{0}
%
\verticaldots{0}{3}
\verticaldots{1}{3}
\verticaldots{2}{3}
\diagonaldots{3}{3}
\verticaldots{4}{3}
\verticaldots{5}{3}
\verticaldots{6}{3}
\diagonaldots{7}{3}
\verticaldots{8}{3}
%
\element{0}{4}{0}        \element{1}{4}{0}        \element{2}{4}{0}
\horizontaldots{3}{4}
\element{4}{4}{-\lambda} \element{5}{4}{0}        \element{6}{4}{0}
\horizontaldots{7}{4}
\element{8}{4}{0}
%
\element{0}{5}{0}        \element{1}{5}{0}        \element{2}{5}{0}
\horizontaldots{3}{5}
\element{4}{5}{0}        \element{5}{5}{1}        \element{6}{5}{-\lambda}
\horizontaldots{7}{5}
\element{8}{5}{0}
%
\element{0}{6}{0}        \element{1}{6}{0}        \element{2}{6}{0}
\horizontaldots{3}{6}
\element{4}{6}{0}        \element{5}{6}{0}        \element{6}{6}{1}
\horizontaldots{7}{6}
\element{8}{6}{0}
%
\verticaldots{0}{7}
\verticaldots{1}{7}
\verticaldots{2}{7}
\diagonaldots{3}{7}
\verticaldots{4}{7}
\verticaldots{5}{7}
\verticaldots{6}{7}
\diagonaldots{7}{7}
\verticaldots{8}{7}
%
\element{0}{8}{0} \element{1}{8}{0} \element{2}{8}{0}
\horizontaldots{3}{8}
\element{4}{8}{0} \element{5}{8}{0} \element{6}{8}{0}
\horizontaldots{7}{8}
\element{8}{8}{1}
\end{tikzpicture}}
%
\;
\right|
=
(-\lambda)^k.
\]
Damit tragen diese Unterdeterminanten insgesamt
\begin{align}
\sum_{k=0}^{n-2}
(-1)^{n+k}
\det \bm{\alpha}_{k+1,n}
&=
\sum_{k=0}^{n-2}
(-1)^{n+k}
(-m_k)
(-\lambda)^k
=
-
(-1)^n
\sum_{k=0}^{n-2}
m_k\lambda^k
\label{buch:komplex:fundamentalsatz:spur:eqn:unterdet}
\end{align}
bei.
Die Minordeterminante $\bm{\alpha}_{n,n}$ für das Element in der rechten
unteren Ecke lässt sich auf analoge Weise bestimmen, sie ist
$(-\lambda)^{n-1}$.
Sie leistet den Beitrag
\[
(-\lambda-m_{n-1})(-\lambda)^{n-1}
=
(-\lambda)^n - m_{n-1}(-\lambda)^{n-1}
=
-
(-1)^n
(\lambda^n + m_{n-1}\lambda^{n-1}).
\]
Zusammen mit \eqref{buch:komplex:fundamentalsatz:spur:eqn:unterdet}
folgt, dass die Determinante von $\bm{\alpha}-\lambda I$ das charakteristische
Polynom
\[
\chi_{\bm{\alpha}}(\lambda)
=
\det (\bm{\alpha}-\lambda I)
=
(-1)^n\biggl(
\lambda^n
+
\sum_{k=0}^{n-1} m_k \lambda^k
\biggr)
=
(-1)^n m(\lambda)
\]
ist.
Die Eigenwerte der Matrix sind also genau die Nullstellen des
ursprünglichen Minimalpolynoms $m(\lambda)$.

%
% Die Spur und die Determinante eines algebraischen Elementes
%
\subsubsection{Die Spur und die Determinante eines algebraischen Elements}
Die Multiplikation mit dem Element $\alpha$ wird in der Basis der
Potenzen von $\alpha$ durch die Matrix $\bm{\alpha}$ 
beschrieben.
Der Matrix $\bm{\alpha}$ kann man einige wohlbekannte Invarianten
zuordnen.

\begin{definition}[Spur und Determinante]
Die \emph{Spur} des algebraischen Elements $\alpha$ ist die Spur der
\index{Spur}%
\index{Determinante}%
Matrix $\bm{\alpha}$, die \emph{Determinante} von $\alpha$ ist die
Determinante der Matrix $\bm{\alpha}$:
\begin{equation}
\operatorname{tr}\alpha
=
\operatorname{tr}\bm{\alpha}
=
-m_{n-1}
\qquad\text{und}\qquad
\det\alpha
=
\det\bm{\alpha}
=
(-1)^nm_0.
\label{buch:komplex:symmetrisch:eqn:spurdet}
\end{equation}
\end{definition}

Die in der Definition enthaltene Aussage, dass Spur und Determinante
von $\alpha$ bereits durch die Koeffizienten des Minimalpolynoms bestimmt
sind, muss erst noch gerechtfertigt werden.
Die Spur und die Determinante einer Matrix sind Eigenschaften,
die nicht von der Wahl einer Basis abhängig sind.
Die Spur und die Determinante der linearen Abbildung
\[
\alpha\colon \mathbb{Q}(\alpha) \to \mathbb{Q}(\alpha) : x \mapsto \alpha x
\]
können also in irgend einer Basis berechnet werden.
Die Wahl von $\mathscr{B}=\{1,\alpha,\dots,\alpha^{n-1}\}$ als Basis
und die daraus resultierende spezielle Form der Matrix $\bm{\alpha}$,
führt auf die Werte \eqref{buch:komplex:symmetrisch:eqn:spurdet}.

\begin{korollar}
Sei $m(X)\in K[X]$ ein irreduzibles Polynom vom Grad $n$ mit Koeffizienten
in $K$.
Seien $\alpha_i$, $i=1,\dots,n$ die Nullstellen von $m(X)$.
Dann ist
\[
\operatorname{tr}\alpha_k
=
\sum_{i=1}^n \alpha_i
=
-m_{n-1}
\qquad\text{und}\qquad
\det \alpha_k
=
\prod_{i=1}^n \alpha_i
=
(-1)^nm_0
\]
für alle $k=1,\dots,n$.
\end{korollar}

\begin{proof}
Die Spur und die Determinante von $\alpha_k$ ist die Spur bzw.~die Determinante
der Matrix $\bm{\alpha_k}$.
Alle Nullstellen haben jedoch die gleiche, aus $m(X)$ konstruierte Matrix
\eqref{buch:komplex:fundamental:eqn:alphaMatrix},
sie haben also alle die gleiche Spur und Determinante, die ausserdem
die Summe bzw.~das Produkt der Eigenwerte dieser Matrix ist.
Das charakteristische Polynom der Matrix ist aber identisch mit dem
Minimalpolynom, die $\alpha_i$ sind somit auch die Eigenwerte.
\index{Eigenwert}%
\end{proof}

\begin{beispiel}
Sei $m(X) = X^3 + X + 1$.
Das Polynom $m(X)$ hat drei verschiedene Nullstellen.
Sei $\alpha$ eine dieser Nullstellen, dann ist
\[
\det \alpha = -1
\qquad\text{und}\qquad
\operatorname{tr}\alpha = 0.
\]
Tatsächlich kann man mit \texttt{octave} die numerischen Werte
\index{octave@\texttt{octave}}%
\[
\renewcommand{\arraycolsep}{2pt}
\begin{array}{rclcl}
\alpha_1 &=& \phantom{-}0.341163901914009 &+& 1.161541399997252i \\
\alpha_2 &=& \phantom{-}0.341163901914009 &-& 1.161541399997252i \\
\alpha_3 &=&           -0.682327803828019 & &
\end{array}
\]
finden.
Die Summe und das Produkt der Nullstellen sind
\[
\renewcommand{\arraycolsep}{2pt}
\begin{array}{rclcl}
\alpha_1+\alpha_2+\alpha_3 &=& -7.771561172376096\cdot 10^{-16} &\approx& -m_2,\\
\alpha_1\alpha_2\alpha_3   &=& -1                               &=&-m_0.
\end{array}
\]
Dies sind bis auf Maschinengenauigkeit die Koeffizienten des
Minimalpolynoms.
\end{beispiel}

%
% Symmetrische Funktionen der Nullstellen
%
\subsubsection{Symmetrische Funktionen der Nullstellen}
Die elementarsymmetrischen Polynome ausgewertet auf den Nullstellen
eines Polynoms mit Koeffizienten in $K$ ergeben Werte in $K$, selbst
dann wenn die Nullstellen selbst nicht in $K$ sind.
Zusammen mit dem
Fundamentalsatz~\ref{buch:komplex:fundamental:satz:fundamentalpolynom}
kann man daraus folgern, dass jedes symmetrische Polynom der Nullstellen
ebenfalls in $K$ liegt.

\begin{satz}
Sei $f(X_1,\dots,X_n)$ ein Polynom mit Koeffizienten in $K$,
das symmetrisch in den Variablen $X_1,\dots,X_n$ ist.
Sei weiter $p(X)\in K[X]$ ein Polynom mit den Nullstellen
$\alpha_1,\dots,\alpha_n$.
Dann ist $f(\alpha_1,\dots,\alpha_n)\in K$.
\end{satz}

\begin{proof}
Nach Satz~\ref{buch:komplex:fundamental:satz:fundamentalpolynom}
kann das Polynom $f(X_1,\dots,X_n)$ als Polynom $g(S_1,\dots,S_n)$
mit Koeffizienten in $K$ geschrieben werden, so dass
\[
f(X_1,\dots,X_n)
=
g(\sigma_1(X_1,\dots,X_n),\dots,\sigma_n(X_1,\dots,X_n))
\]
gilt.
Einsetzen der Nullstellen ergibt
\[
f(\alpha_1,\dots,\alpha_n)
=
g(
\underbrace{\sigma_1(\alpha_1,\dots,\alpha_n)}_{\displaystyle\in K}
,\dots,
\underbrace{\sigma_n(\alpha_1,\dots,\alpha_n)}_{\displaystyle\in K}
)
\in K.
\]
Es folgt, dass der Wert $f(\alpha_1,\dots,\alpha_n)$ in $K$ liegt.
\end{proof}

Wenn man zeigen kann, dass eine rationale Funktion von $X_1,\dots,X_n$
genau dann symmetrisch ist, wenn Zähler und Nenner symmetrisch sind,
dann folgt auch, dass die Werte rationaler, symmetrischer Funktionen
der Nullstellen im Koeffizientenkörper liegen.

\begin{satz}
Sei $f(X_1,\dots,X_n)$ eine rationale Funktion mit Koeffizienten
in $K$, die symmetrisch ist unter Vertauschungen der Argumente.
Dann kann $f$ geschrieben werden als Bruch
\[
f(X_1,\dots,X_n)
=
\frac{u(X_1,\dots,X_n)}{v(X_1,\dots,X_n)},
\]
wobei sowohl $u$ als auch $v$ symmetrische Polynome mit Koeffizienten
in $K$ sind.
Insbesondere ist $f(\alpha_1,\dots,\alpha_n)$ auch eine rationale Funktion
der Werte $\sigma_k(\alpha_1,\dots,\alpha_n)\in K$ mit Koeffizienten in $K$.
\end{satz}

\begin{proof}
Sei 
\[
f(X_1,\dots,X_n)
=
\frac{a(X_1,\dots,X_n)}{b(X_1,\dots,X_n)}
\]
mit Polynomen $a,b\in K[X_1,\dots,X_n]$.
Die rationale Funktion $f$ ist symmetrisch.
Insbesondere folgt
\begin{align*}
f(X_1,\dots,X_n)
&=
\frac{1}{n!}
\sum_{\pi\in S_n}
\frac{a^\pi(X_1,\dots,X_n)}{b^\pi(X_1,\dots,X_n)}.
\intertext{Durch Zusammenfassen entsteht ein Bruch der Form}
&=
\frac{
u(X_1,\dots,X_n)
}{
v(X_1,\dots,X_n)
},
\qquad
\text{mit }
v(X_1,\dots,X_n)
=
n!\prod_{\pi\in S_n} b^\pi(X_1,\dots,X_n).
\end{align*}
Der Nenner $v(X_1,\dots,X_n)$ ist offensichtlich ein symmetrisches
Polynom.
Da auch $f$ symmetrisch ist, folgt, dass das Polynom
\[
u(X_1,\dots,X_n)
=
f(X_1,\dots,X_n)\cdot v(X_1,\dots,X_n)
\]
ebenfalls symmetrisch ist.
Damit ist die Existenz der symmetrischen Polynome $u$ und $v$
gezeigt.
\end{proof}

\begin{korollar}
Sei $p(X)\in K[X]$ ein Polynom mit den Nullstellen
$\alpha_1,\dots,\alpha_n$ und
$f(X_1,\dots,X_n)$ eine rationale Funktion mit Koeffizienten
in $K$, die symmetrisch ist unter Vertauschungen der Argumente.
Dann ist $f(\alpha_1,\dots,\alpha_n)$ auch eine rationale Funktion
der Werte $\sigma_k(\alpha_1,\dots,\alpha_n)\in K$ mit Koeffizienten in $K$.
\end{korollar}

Das Korollar bedeutet, dass man die Nullstellen eines Polynoms 
nicht kennen muss, wenn man die Werte einer symmetrischen rationalen
Funktion der Nullstellen berechnen will, weil diese Werte bereits vollständig
durch die Koeffizienten des Polynoms bestimmt sind.

